\documentclass[10pt,a4paper]{amsart}
\usepackage[margin=19mm]{geometry}
\usepackage{amsmath,amssymb,mathtools}
\usepackage[scr=rsfs,cal=euler]{mathalfa}
\usepackage{xfrac}
\usepackage[unicode,hidelinks,bookmarks=false]{hyperref}
\newcommand{\ie}{i.e.\ }
\newcommand{\cf}{cf.\ }
\newcommand{\resp}{respectively\ }
\newcommand{\Subsection}{\subsection}
\providecommand{\nobreakdash}{\nobreak\hbox{-}}
\setlength{\emergencystretch}{2em}
\multlinegap0pt
\usepackage{bm}
\usepackage{bbm}
\usepackage{dsfont}
\usepackage{xspace}
\usepackage{cancel}
\usepackage{mathtools}
\usepackage{amsthm}
\makeatletter
\@ifundefined{theorem}{\newtheorem{theorem}{Theorem}}{}
\@ifundefined{lemma}{\newtheorem{lemma}[theorem]{Lemma}}{}
\makeatother
\title[Impartial Combinatorial Games and the Nuclear Escalation Lad]{Impartial Combinatorial Games and the Nuclear Escalation Ladder}
\author{Arnav Garg}
\date{Source: arXiv:2606.29015v1 (CC BY 4.0)}
\begin{document}
\maketitle

\noindent\emph{Source and licence.} Impartial Combinatorial Games and the Nuclear Escalation Ladder. Version arXiv:2606.29015v1 (CC BY 4.0), distributed under the Creative Commons Attribution 4.0 International licence, as recorded in the arXiv licence metadata for this exact version and shown on the abstract page https://arxiv.org/abs/2606.29015v1. Full rights notice: licenses/arxiv-2606.29015-CC-BY-4.0.txt.\par\smallskip

\noindent\textit{Editorial scope.} Counted unit: the Theorem whose source label is \texttt{thm:B} in the article (source0.tex lines 470--483, source label \texttt{thm:B}), delivered together with the article's own complete original proof of it (source0.tex lines 485--525, one \texttt{proof} environment of 41 lines). No mathematics is written, repaired or re-derived: statement and proof are the article's own text line for line. Required context retained uncounted: lem:grundy (lines 437--446). Every definition, display and label of those lines is retained unchanged. Omitted from this package and not redistributed: 1--436, 447--469, 484--484, 526--879. Nothing inside a retained excerpt is omitted or altered: every retained body is delivered exactly as the article's lines are written, so no omission receipt of any kind accompanies this package. Citations: the retained excerpts carry no citation command at all, so no bibliography entry of the article is transcribed and no reference is redistributed. Reference adaptation: none was needed and none was made; every reference inside the delivered unit resolves inside it, so no unresolved reference or citation is produced. Printed slips kept literal: no printed word, symbol, hypothesis or number of the counted statement or of its proof was altered. Layout and class: the article's own class is \texttt{article} with packages including lmodern, fontenc, inputenc, amsmath, amssymb, amsthm, microtype, natbib, geometry, hyperref; the wrapper is the lane's pinned \texttt{amsart} editorial wrapper at 10pt on A4, which restates the theorem-like environments the retained text needs and adds the article's own macro definitions that the retained text needs (none), with extra wrapper packages (bm, bbm, dsfont, xspace, cancel, mathtools). The delivered numbering is the unit's own. Assets: no retained span contains a figure environment, a graphics-inclusion command, a PDF-page inclusion, a table or a TikZ picture; no image file is copied into this package, the package declares no asset dependency and no image file is redistributed with it. Licence: version arXiv:2606.29015v1 of this article is distributed under the Creative Commons Attribution 4.0 International licence, as recorded in the arXiv licence metadata for this exact version and shown on the abstract page https://arxiv.org/abs/2606.29015v1. A full rights notice is delivered at licenses/arxiv-2606.29015-CC-BY-4.0.txt. Characters: every retained excerpt is pure ASCII, so the delivered engine, XeLaTeX with its default Latin Modern text font, reports no missing character.
\begin{lemma}[Component Grundy values and reachability]\label{lem:grundy}
In normal play, the single-theater subtraction game with escalation set
$S_i=\{1,\dots,m_i\}$ has Grundy value
\[
  G_i(r) \;=\; r \bmod (m_i+1), \qquad r \ge 0 .
\]
Moreover, for every state $r\ge 0$ and every value
$v \in \{0,1,\dots,G_i(r)-1\}$ there is a legal move from $r$ to some
state of Grundy value $v$, and no legal move preserves the Grundy value.
\end{lemma}

\begin{theorem}[Multi-theater Nim-sum stability]\label{thm:B}
Consider $k$ independent theaters, theater $i$ in state $r_i$ with
escalation set $S_i=\{1,\dots,m_i\}$. On each turn the player to move
selects exactly one theater $i$ and subtracts some $d\in S_i$ with
$d\le r_i$. Under normal play, the joint state $(r_1,\dots,r_k)$ is a
$\mathcal{P}$-position if and only if
\[
  \bigoplus_{i=1}^{k} G_i(r_i)
  \;=\;
  \bigoplus_{i=1}^{k}\big(r_i \bmod (m_i+1)\big)
  \;=\; 0,
\]
where $\oplus$ denotes the bitwise exclusive-or (Nim-sum).
\end{theorem}

\begin{proof}
Write $g_i := G_i(r_i)$ and $s := \bigoplus_{i=1}^{k} g_i$. The terminal
state is $(0,\dots,0)$, which has no moves and is therefore a
$\mathcal{P}$-position; its Nim-sum is $s=0$. We prove by induction on the
total $\sum_i r_i$ that $s=0$ characterizes $\mathcal{P}$-positions,
establishing the two standard claims.

\emph{Claim 1 (from $s=0$, every move yields $s\neq 0$).}
A move changes exactly one component, say theater $j$, replacing $r_j$ by
$r_j-d$ and hence $g_j$ by $g_j':=G_j(r_j-d)$. By Lemma~\ref{lem:grundy} no
move preserves a component's Grundy value, so $g_j'\neq g_j$. The new
Nim-sum is
\[
  s' \;=\; s \oplus g_j \oplus g_j' \;=\; 0 \oplus g_j \oplus g_j'
       \;=\; g_j \oplus g_j' \;\neq\; 0,
\]
since $g_j\neq g_j'$ implies $g_j\oplus g_j'\neq 0$. Thus every option of a
state with $s=0$ has nonzero Nim-sum, and by the induction hypothesis is
an $\mathcal{N}$-position. A state all of whose options are
$\mathcal{N}$ is a $\mathcal{P}$-position.

\emph{Claim 2 (from $s\neq 0$, some move yields $s=0$).}
Let $b$ be the position of the highest set bit of $s$. Some component $g_j$
has bit $b$ set, because the bits of $s$ are the parity of the components'
bits. Put $g_j' := g_j \oplus s$. Flipping bit $b$ of $g_j$ from $1$ to $0$
and possibly altering lower bits strictly decreases the value, so
$g_j' < g_j$. By the reachability part of Lemma~\ref{lem:grundy}, theater
$j$ admits a legal move from $g_j$ to a state of Grundy value $g_j'$.
After that move the Nim-sum is
\[
  s' \;=\; s \oplus g_j \oplus g_j'
       \;=\; s \oplus g_j \oplus (g_j \oplus s) \;=\; 0 .
\]
By the induction hypothesis that option is a $\mathcal{P}$-position, so the
original state, having a $\mathcal{P}$-option, is an
$\mathcal{N}$-position.

Claims 1 and 2 together give: $s=0$ if and only if the state is a
$\mathcal{P}$-position. Substituting $g_i=r_i\bmod(m_i+1)$ from
Lemma~\ref{lem:grundy} yields the stated congruence form.
\end{proof}

\end{document}
